\section{Zadanie 12}
Czy jeśli dla dwóch zmiennych losowych X, Y mamy Cov(X, Y ) = 0 to
X, Y są niezależne?

\subsection*{Rozwiązanie}
Nie jest tak, można łatwo znaleźć kontrprzykład. Załóżmy, że z równym prawdopodobieństwem losujemy jedną z liczb \(\{-1,0,1\}\). Wówczas niech X - zmienna losowa oznaczająca wylosowaną liczbę. Zauważmy, że \(E[X] = (-1)\cdot P(X=-1) + 1\cdot P(X=1) + 0\cdot P(X = 0)= \frac{1}{3}\cdot (-1) + \frac{1}{3} \cdot 1 = 0\). Analogicznie \(E[X^3] = 0\). Teraz niech zmienna losowa Y = \(X^2\) oznacza kwadrat wylosowanej liczby. Wtedy \(E[Y] = \frac{1}{3}\cdot 1^2 + \frac{1}{3} \cdot (-1)^2  + \frac{1}{3} \cdot 0^2= \frac{2}{3}\). Oczywiście zmienne X i Y nie są niezależne, np. \(P(X=1 \And Y=0) = 0 \neq \frac{1}{9} =  P(X = 1) \cdot P(Y = 0)\), a mimo to:
\[
	Cov(X,Y) = E[XY] - E[X]E[Y] = E[X^3] - 0 = 0
\]


\subsection*{Odpowiedź}
\[
	\boxed{\text{Nie}}
\]
